\section{Research aims and questions}

The goal of this project is to address some gaps in the current AI alignment landscape by a proof-of-concept end-to-end project. It will attempt to align a large language model to a set of human values collected from the public and to develop an evaluation method that can better measure alignment success on these values. To do this I will use value survey data to do alignment training and evaluate using behavioral assessments of the LLM by the public. Specifically, the project will set out to answer a series of research questions.

\begin{enumerate}[label=\textbf{RQ\arabic*:}, ref=RQ\arabic*, leftmargin=*, itemsep=1em]
  \item \textbf{Alignment feasibility.} Do simple post-training methods (e.g  SFT) work to appreciably align a large language model to a set of human values? \label{rq1}

  \item \textbf{Evaluation.} What method can be used to evaluate the alignment of a large language model to a set of representatively collected human values, in a way that remains democratically grounded? \label{rq2}
\end{enumerate}

These research questions, if addressed, would provide useful insights into the feasibility of aligning large language models to the public in a democratically grounded way\footnote{These research questions are supported by efforts like \cite{anwarFoundationalChallengesAssuring2024}}. This is of particular importance within the context of Public, Sovereign and Democratic AI \cite{daleSovereignAI20252025,nfeltSovereignRemediesAI2025,publicaiPublicAIWhite2024,ovadyaDemocraticAIPossible2025} movements that are pushing for more public participation, oversight, and accountability in AI systems.

This project sets out with some hypotheses around these research questions.


\begin{enumerate}[label=\textbf{H\arabic*:}]
    \item A small ($<$40B parameters) \term{open-weights} LLM has the ability to understand the survey questions and model a respondents values when given its responses. \\
    \textit{Justification:} Small LLMs have been shown to respond positively to fine tuning with value/opinion survey data with the result of better prediction of survey responses \cite{suhLanguageModelFineTuning2025}.
    
    \item Fine-tuning a base LLM by simply telling it which answers to a survey question is \scare{correct} will change its behavior in downstream tasks.\\
    \textit{Justification:} Recent work \cite{nieSurveytoBehaviorDownstreamAlignment2025} has suggested that fine tuning methods similar to this has produced detectable behavioral differences in the LLM in sensitive situations.
    \item The questions and responses in the WVS are useful to predict behavior, as tested by asking respondents how they would act in certain situations.\\
    \textit{Justification:} The foundations of these survey methods is that they allow us to understand humans values, many years of research has suggested that values are good predictors of behaviour \cite{schwartzRefiningTheoryBasic2012, rokeachNatureHumanValues1973}. 
    \item The public will rate models that are fine-tuned to their survey responses as most aligned to them compared with models fine-tuned to other clusters or the baseline model.\\
    \textit{Justification:} Individuals are shown to prefer agents that they perceive share common values with them which is called similarity-attraction \cite{byrneAttractionParadigm1971}.
\end{enumerate}


\section{Alignment methodology and evaluation plan}

To address these research questions, I will need to carry out a sequence of steps from ingesting a pre-trained model to outputting an aligned model and evaluation results.

\subsection{Defining the alignment target}

For pragmatic reasons, I will use existing survey data from the world  to define a set of human values to align towards. This will provide me with a dataset of 300 questions answered by about 1,000 respondents in multiple countries. Either each country will be treated as a separate cluster or I will cluster respondents based on their answers to define value clusters. This will allow me to explore pluralistic alignment targets.

I will define two sorts of alignment targets. The first will be a set of global values that represent the values of the population as a whole, this can be derived by taking the mode of the entire response along with a variation metric. The second will be pluralistic targets that get defined as the mode response within each cluster of respondents (one cluster for each country). This will allow me to explore the differences between different value alignment targets.

\subsection{Aligning to the targets}

Recent work \cite{nieSurveytoBehaviorDownstreamAlignment2025} suggests that simple SFT may have appreciable alignment effects. Therefore, I will attempt to use SFT by getting the model to answer the survey questions correctly, in Figure \ref{fig:wvs-simple-sft} one can see the basic outline. If this proves insufficient, I will explore more complex methods such as KTO with binary feedback (correct survey response, incorrect survey response). These methods will be used to align a base LLM to the defined alignment targets from step 1.

Given there are only about 290 survey questions, I will need to augment the data to ensure there are enough training examples. I will do this by using different templates that pose the same question in different ways. Importantly the actual content of the question will remain the same, only the system prompt and structure of getting the LLM to answer will change. By having a set of different templates it is plausible that the 290 initial questions can be augmented to create a training set of about 1500 examples.

\begin{figure}[h]
\centering
\begin{tikzpicture}[
    every node/.style={draw, rounded corners, align=left, font=\small},
    node distance=1.4cm,
    arrow/.style={->, thick}
]

\node (prompt) {
\textbf{Model input (Prompt)}\\[2pt]
Instructions: On a scale from 1 to 10 where 1 means 'not at all important' and 10 means 'absolutely important',\\rate how important it is:\\``for you to live in a country that is governed democratically?''};

\node (model) [below=0.3cm of prompt] {
\textbf{Large Language Model}
};

\node (output) [below of= model] {
\textbf{Model output}\\[2pt]
6
};

\node (real) [right=1cm of output] {
\textbf{Response from WVS respondents}\\[2pt]
7
};

\node (feedback) [below=1cm of output] {
Loss = $\mathcal{L}(pred, real)$
};

\draw[arrow] (prompt) -- (model);
\draw[arrow] (model) -- (output);
\draw[arrow] (output) -- (feedback);
\draw[arrow] (real) -- (feedback);
\draw[arrow, dashed] 
    (feedback.west) 
    to[bend left=25] 
    node[left, font=\small]{Update model to increase the\\ likelihood of generating true response}
    (model.west);

\end{tikzpicture}
\caption{Simplified illustrative example using WVS survey data. The model is prompted with a survey question, produces a Likert-scale response, and receives a supervision signal based on its distance from the observed WVS response.}
\label{fig:wvs  -simple-sft}

\end{figure}

\subsection{Evaluating the success of the alignment}

To evaluate the success of the alignment, I will need to develop an evaluation method that can measure how well the AI agent is aligned to the defined human values. There will be broadly two tests. One is automated evaluation using held-out survey questions. The other is a public consultation process where I will collect feedback on how much the person feels the model acted in alignment with their values. This will require significant design to ensure that the consultation is representative and captures the necessary feedback.

The first evaluation method will involve holding out a set of survey questions from the training data. The test split will be done by splitting out questions that are semantically similar to each other and holding out different question templates. This will ensure that the model has at least some ability to generalize to unseen questions.

The second evaluation method will involve getting the model to engage in various simulations of real world tasks. Then I will recruit participants from the various clusters defined in step 1 to rank which model actions best align with their values in these scenarios. This will provide a human-grounded evaluation of alignment success.

These evaluations combined will provide a robust measure of how well the model has understood the human values as well as how it actually affects its behavior in more realistic settings.

\section{Rationale for the project}

Current methods for training large language models involve scraping the Internet for large amounts of text data, creating powerful next-token prediction models. To turn these into something useful for end users (e.g. chatbots or assistants), the model is then aligned to behave in desired ways, typically by contracting human labelers/annotators/demonstrators to provide training data. This has made modern LLMs highly useful.
Problematically, current alignment methods are opaque and the alignment target is often specified privately by a small set of individuals (e.g. the people writing instructions for annotators). This requires trust from the public that the model is aligned with their values. Even if the values of the creators were initially acceptable, they are usually set within for-profit organizations with incentives that can distort alignment (e.g. maximizing engagement or revenue). Therefore we cannot assume these models are aligned to the public good.

There are multiple challenges, however, with aligning models to the public good, or more formally maximizing median human flourishing.
First is how to define what the public good is, that is what would the alignment target even look like? Avenues to explore are; lists of principles (like constitutional AI \cite{baiConstitutionalAIHarmlessness2022}), value sets from a taxonomy of human values \cite{schwartzUniversalsContentStructure1992,rokeachNatureHumanValues1973}, fine grained preference datasets (like that collected for RLHF \cite{ouyangTrainingLanguageModels2022} or similar), or more complex structures like Moral graphs \cite{klingefjordWhatAreHuman2024}.
Second is that even if we understood what an alignment target should look like, it is non-obvious how to collect data that represents it in a way that is representative of the public.
Third is that once this alignment structure is defined and data collected to turn it into a valid instance of an alignment target it is an open question how best to align a subsequent model towards this target. This becomes particularly challenging given that the alignment target could be a pluralistic target and thus require fair distinct treatment of different subgroups.
Finally, even if we had a method to align a model to this target, it is unclear how we would evaluate success. Current methods are generally short term and often take the form of pairwise comparisons or held-out test sets of expected responses. Notably, they do not usually provide out-of-distribution evaluation or capture multi-step interactions and downstream outcomes, and they are often designed and performed by the same select few individuals who designed the alignment target in the first place.

This project aims to address these challenges by providing a demonstration that we could use readily available value survey data to define an alignment target. Then train a model to align to this target using simple methods. Finally, evaluation of alignment in out-of-distribution scenarios can be done by judging the model's behavior in realistic scenarios by members of the public.


\section{Contributions}

The project overall will contribute by providing a proof-of-concept alignment procedure to the NZ public. Within that it will contribute:
\begin{itemize}
  \item A method to define alignment targets and a training dataset from survey data that captures pluralistic human values in a democratically defensible manner.
  \item A novel [application of]\footnote{Unsure if creating my own or using someone else's dataset} evaluation dataset to test a models behavior in realistic scenarios to be judged by human participants from various value clusters.
  \item A demonstration of a lightweight training procedure on open weight LLMs that has democratically fair evaluation across different sub groups.\footnote{Important to note I don't know how successful this will be yet. Ideally I could say that it achieves high alignment scores but this is unknown at this stage.}
\end{itemize}

\noindent All of the work will be open-sourced, documented and freely available.

\section{Resources and feasibility}

This project will take just over 4 months to complete and require a few thousand dollars of compute and participant compensation.

\subsection{Proposed timeline}

Figure \ref{fig:timeline} shows the proposed timeline for a 18-week project.

\begin{figure}[h!]
\centering
\begin{ganttchart}[
  vgrid, hgrid,
  x unit=0.5cm
]{1}{18}

  % --- Month labels ---
  \gantttitle{Feb}{3}
  \gantttitle{Mar}{4}
  \gantttitle{Apr}{4}
  \gantttitle{May}{4}
  \gantttitle{Jun}{3} \\

  % --- Week numbers ---
  \gantttitlelist{1,...,18}{1} \\

  % --- Work packages ---
  \ganttgroup{Survey Data Analysis}{1}{2} \\
  \ganttbar{Collect and organise data}{1}{1} \\
  \ganttbar{Format into dataset}{1}{2} \\

  \ganttgroup{Aligning Model}{2}{8} \\
  \ganttbar{Select Base Model}{2}{2} \\
  \ganttbar{Pilot fine tune to test pipeline}{3}{4} \\
  \ganttbar{Full fine tuning runs}{7}{8} \\
  
  \ganttgroup{Evaluating Alignment}{5}{17}\\
  \ganttbar{Create Eval Method}{5}{7} \\
  \ganttbar{Initial Evaluation}{9}{10} \\
  \ganttbar{Final Evaluation}{15}{16} \\

  \ganttgroup{Public consultation}{7}{14} \\
  \ganttbar{Ethics approval application}{7}{8} \\
  \ganttbar{Design}{11}{13} \\
  \ganttbar{Collect Responses}{13}{14} \\

  \ganttgroup{Writing report}{8}{18} \\
  \ganttbar{Writeup alignment process}{8}{9} \\
  \ganttbar{Writeup Eval method}{9}{10} \\
  \ganttbar{Writeup public consultation}{14}{14} \\
  \ganttbar{Complete writeup}{16}{18} 
\end{ganttchart}

\caption{Project timeline in weeks (only roughly aligned to months). This assumes a start time of 9th of February 2026.}
\label{fig:timeline}
\end{figure}

\subsection{Computational resources}

To train and evaluate the model, hardware will be needed. If I can train models with about 30 billion parameters, I will be able to get near-SOTA performance \footnote{Based on the \href{https://artificialanalysis.ai/leaderboards}{Artificial Analysis leaderboards}: 30B parameters achieve about 60\% of the performance of frontier models.}. Using modern quantization methods \cite{dettmersQLoRAEfficientFinetuning2023,dettmersLLMint88bitMatrix2022} and parameter-efficient fine-tuning methods \cite{huLoRALowRankAdaptation2021,dettmersQLoRAEfficientFinetuning2023}, I should be able to fine-tune these models using around 80GB of VRAM (h100, h200 etc). Therefore, I estimate needing about 240 USD of GPU time to carry out the fine-tuning and evaluation if hardware is rented in the cloud \footnote{Each fine-tuning run will take in the realm of 3--6 hours depending on the size of the model, amount of data used, alignment methods used, and GPU hardware. About 80 hours (5 hours per fine-tune * 3 different methods * 5 different value sets + 5 hours of evaluation) of GPU time will be needed at a cost of 2--3 USD for a single hour on a high-end GPU (e.g. H200) \cite{vastai2026}}. Alternatively the university has some GPU resources that could be used and might be sufficient.

\subsection{Survey data}

The initial alignment target will be defined using existing survey data from the World Values Survey (WVS) \cite{WVS2020} survey data from across the globe with about 1,000 respondents in each country. It is openly accessible for research purposes \footnote{Initial plans of using NZAVS data was not possible due to failure to obtain access}.

\subsection{Public consultation responses}

To create a rigorous evaluation of the alignment process a public consultation process will be needed. This will require recruiting participants from various clusters of the NZ population to provide feedback. Each participants will rate how much an AI agents actions (i.e responses) align with their values. I would expect to need at least 50 people from various countries to complete about a 30 minutes task (including instructions). The complete design of this is left to the research phase.

\section{Ethical considerations}

The interaction of this project with human data and human participants raises ethical concerns. Furthermore as AI has the potential for strong negative real-world impact there are considerations around development and training of these models.

\begin{itemize}
  \item The use of human survey data (WVS) to construct alignment targets. \\
  \textit{mitigation:} Only aggregate-level data will be used, and all data custodians' access and usage conditions will be followed.
  \item Interaction with human participants during the public consultation evaluation.\\
  \textit{mitigation:} Ethics approval will be sought prior to recruitment, and informed consent will be obtained from all participants.
  \item Concerns of dual-use and capability increases from alignment work. \\
  \textit{mitigation:} The project will use existing open-weight models (not frontier-scale training) and focus on value alignment rather than capability enhancement. Concerns of emergent misalignment will be addressed by only training to ever do \scare{positive} things.
\end{itemize}